{%
\def\insertshortauthor{吕昀键}
\def\insertshorttitle{07.20\textasciitilde07.26 工作进展}
\def\insertshortdate{2026.07.27}
\begin{frame}[t]{跨 GPU Longer-MQAR: 优势集中在长度外推}
  \centering
  \vspace{0.06cm}
  \begin{minipage}{0.94\textwidth}
  \begin{columns}[T,onlytextwidth]
    \begin{column}{0.71\linewidth}
      \makebox[\linewidth][c]{%
        \includegraphics[
          width=1.17\linewidth,
          height=5.72cm,
          keepaspectratio
        ]{20260720-0726/p18/longer-mqar-accuracy-last.png}%
      }
    \end{column}

    \begin{column}{0.26\linewidth}
      \thinsectionbar{实验覆盖}
      \vskip0.05cm
      \begin{beamercolorbox}[sep=3pt,rounded=true,wd=\linewidth]{block body}
        \fontsize{6.5}{7.2}\selectfont
        2 GPU $\times$ 2 模型 $\times$ 3 seeds, \textcolor{scured}{\bfseries 12/12} 个正式 4 epoch run.
      \end{beamercolorbox}

      \vskip0.07cm
      \thinsectionbar{质量结果}
      \vskip0.05cm
      \begin{beamercolorbox}[sep=3pt,rounded=true,wd=\linewidth]{block body}
        \fontsize{6.5}{7.2}\selectfont
        端点 1024$\times$256: Flash 不领先.\par
        \vskip0.03cm
        4 个外推切片: Flash 平均准确率\textcolor{scured}{\bfseries 均高于 GDN}.\par
        \vskip0.03cm
        RTX 3090: 4/4 切片均为 \textbf{3/3 seeds} 领先.
      \end{beamercolorbox}

      \vskip0.07cm
      \thinsectionbar{跨 GPU 边界}
      \vskip0.05cm
      \begin{beamercolorbox}[sep=3pt,rounded=true,wd=\linewidth]{block body}
        \fontsize{6.5}{7.2}\selectfont
        GDN 同 seed 稳定; Flash 主要波动来自 \textbf{seed124}.\par
        \vskip0.03cm
        每张 GPU 按 \textbf{$n=3$} 统计, 不合并为 $n=6$.
      \end{beamercolorbox}
    \end{column}
  \end{columns}

  \vskip0.05cm
  \begin{beamercolorbox}[sep=3pt,rounded=true,wd=\linewidth]{block body}
    \centering\small
    \textcolor{scured}{\bfseries 结论: Flash 的优势来自训练长度之外的关联记忆保持, 但训练轨迹对 seed 和 GPU 数值路径更敏感.}
  \end{beamercolorbox}
  \end{minipage}
\end{frame}
}
